\documentclass[11pt]{article}
\usepackage{amsmath,amssymb,amsthm,geometry,hyperref}
\geometry{margin=25mm}
\newtheorem{theorem}{Theorem}
\newcommand{\Tr}{\operatorname{Tr}}
\begin{document}
\title{A spin-three inertia bound and the fourth-order maximum at eighteen qubits}
\author{Exact catalogue calculation}
\date{3 November 2026}
\maketitle
Write $K(N,t)$ for the largest dimension of a complex-linear subspace of spin $N/2$
every normalized vector of which is $t$-anticoherent.
\begin{theorem}
For every integer $N\ge16$ and $t\ge4$,
\[
 \boxed{7K(N,t)\le N-5.}
\]
In particular, a two-dimensional fourth-order AC subspace on eighteen
qubits does not exist. The existing exact state attains dimension one:
\[
 \boxed{K(18,4)=1.}
\]
\end{theorem}
The proof applies to arbitrary complex subspaces. It makes no assumption
about Dicke support, constellation symmetry or reality of amplitudes.

\section*{1. Anticoherence gives a zero compression}
Let $P$ project onto a $k$-dimensional fourth-order AC subspace of $V_j$,
where $j=N/2$. Polarization of the AC conditions gives
\[
 PT_{\ell q}P=0\qquad(1\le\ell\le4)
\]
for all irreducible spin-tensor components. Every ordered polynomial in
$J_x,J_y,J_z$ of degree at most four has only tensor ranks zero through
four. Therefore, for every such operator $A$,
\[
 PAP=\frac{\Tr A}{2j+1}P. \tag{1}
\]
Introduce an auxiliary spin three, with operators $S_x,S_y,S_z$ on $V_3$,
and put
\[
 X=J_x\otimes S_x+J_y\otimes S_y+J_z\otimes S_z.
\]
If $W$ is a polynomial of degree at most four in $X$, (1) applies to each
physical-spin coefficient of $W$. Thus
\[
 (P\otimes I_7)W(P\otimes I_7)
 =\frac{P}{2j+1}\otimes\Tr_{V_j}W. \tag{2}
\]
We will construct $W$ with zero partial trace and only $N-5$ positive
eigenvalues, counted with multiplicity.

\section*{2. An explicit quartic with the required inertia}
For $j\ge8$, define
\[
 F_j(x)=(x+2j+5)(x+6)(x-j+5)(x-3j),
\]
\[
 c_j=2j(4j^3-36j^2+47j-75),\qquad W=F_j(X)-c_j I.
\]
The Clebsch--Gordan decomposition
\[
 V_j\otimes V_3=\bigoplus_{p=0}^{6}V_{j-3+p}
\]
gives the eigenvalue of $X$ at total spin $L$ as
\[
 x_L=\frac{L(L+1)-j(j+1)-12}{2}.
\]
The complete spectrum is consequently
\[
\begin{array}{c|c|c}
p&x_{j-3+p}&F_j(x_{j-3+p})\\\hline
0&-3j-3&18(j-2)(j-1)(2j-1)(2j+1)\\
1&-2j-5&0\\
2&-j-6&-2j(j-1)(2j+1)(2j+3)\\
3&-6&0\\
4&j-5&0\\
5&2j-3&-2(j+2)(j+3)(2j+1)(2j+3)\\
6&3j&0
\end{array}
\]
with multiplicities $2j-5+2p$. Summing these seven weighted values gives
\[
 \Tr F_j(X)=7(2j+1)c_j. \tag{3}
\]
Since $F_j(X)$ is invariant under the joint rotation representation,
its partial trace over $V_j$ commutes with the irreducible spin-three
representation. It is scalar; (3) determines that scalar:
\[
 \Tr_{V_j}F_j(X)=(2j+1)c_j I_7,
 \qquad \Tr_{V_j}W=0. \tag{4}
\]
For $h=j-8\ge0$,
\[
 c_j=2j(4h^3+60h^2+239h+45)>0.
\]
The table shows that the six sectors $p=1,\ldots,6$ have strictly negative
eigenvalues of $W$. Its remaining eigenvalue must be strictly positive:
$\Tr W=0$ and the other sectors have negative trace. Equivalently, it is
\[
 F_j(-3j-3)-c_j
 =4(j+1)(16j^3-52j^2+60j-9)>0.
\]
Thus $W$ is nonsingular and has positive and negative indices
\[
 n_+(W)=2j-5=N-5,\qquad n_-(W)=6(N+2). \tag{5}
\]

\section*{3. Inertia bounds the whole AC subspace}
Equations (2) and (4) show that
$\mathcal U=\operatorname{Ran}P\otimes V_3$, of dimension $7k$, is totally
isotropic for $W$: its entire compressed matrix is zero.
For a nonsingular Hermitian form, projection of a totally isotropic
subspace onto its positive spectral subspace is injective. Indeed, a
nonzero vector with zero positive projection would belong to the strictly
negative spectral subspace and have negative quadratic form, a
contradiction. Therefore
\[
 7k=\dim\mathcal U\le n_+(W)=N-5.
\]
Every $t$-AC subspace with $t\ge4$ is also fourth-order AC, proving the
stated general bound.

\section*{4. Specialization to eighteen qubits}
For $j=9$ one obtains the particularly simple operator
\[
 \begin{aligned}
 W&=(X+23)(X+6)(X-4)(X-27)-6264 I\\
  &=X^4-2X^3-653X^2-1146X+8640 I.
 \end{aligned}
\]
Its exact spectrum is
\[
\begin{array}{c|rrrrrrr}
L&6&7&8&9&10&11&12\\\hline
2L+1&13&15&17&19&21&23&25\\
W|_{V_L}&319320&-6264&-63720&-6264&-6264&-111600&-6264
\end{array}
\]
so its positive index is thirteen. A two-dimensional AC4 subspace would
produce a fourteen-dimensional totally isotropic subspace, which is
impossible. Thus $K(18,4)\le\lfloor13/7\rfloor=1$.

The supplied Julia source and AC-state catalogue already contain
\[
 |\psi\rangle=\frac{\sqrt{14}(|D_{18}^{(0)}\rangle+|D_{18}^{(18)}\rangle)
 +i\sqrt{51}(|D_{18}^{(5)}\rangle+|D_{18}^{(13)}\rangle)}{\sqrt{130}}.
\]
Its probabilities are $(14,51,51,14)/130$, so it is normalized and
reflection symmetry cancels odd diagonal moments. Moreover
\[
 \langle J_z^2\rangle=30=\frac{9\cdot10}{3},\qquad
 \langle J_z^4\rangle=1614
 =\frac{9\cdot10(3\cdot9^2+3\cdot9-1)}{15}.
\]
Every support gap is at least five, so all nonzero magnetic shifts through
rank four vanish. Hence this state is fourth-order AC; in fact the known
phase cancellation makes it fifth-order AC. Its one-dimensional span
attains the upper bound.

\paragraph{Independent exact verification.}
\texttt{verify\_upper.py} reconstructs the physical spin-nine and
spin-three matrices in an integer magnetic basis. It checks weighted
Hermiticity, all 25 magnetic-block characteristic polynomials, all 49
entries of the partial trace, and the complete inertia using exact
arithmetic. No numerical eigenvalues enter the proof. The general trace
identity is also checked by exact polynomial interpolation with its
explicit degree bound. \texttt{upper-verification.json} records the result.
The lower witness is checked independently against its exact Dicke
coefficients and reduced states in the accompanying lower verification.

At nineteen qubits this theorem alone gives $K(19,4)\le2$. The separate exact four-copy certificate in \texttt{../n19-t4/exclusion/proof.tex} now proves $K(19,4)=1$.
\end{document}
